%%%%%%%%%%%%%%%%%%%%%%%%%%%%%%%%%%%%%%%%%%%%%%%%%%%%%%%%%%%%%%%%%%%%%%%%%%%%%%%%%%%%%%%%%%%%%%%%%%%%%%%%%%%%%%%%%%%%%%%%%%%%%%%%%%%%%%%
%%| !!! TEMPLATE INSPIRED BY https://github.com                  !!! |%%
%%| !!! HsKa_old.jpg CAN BE FOUND HERE: https://de.wikipedia.org   !!! |%%
%%| !!! HsKa_new.png WAS CREATED BY USAGE OF IMAGERY FROM: https://www.h-ka.de       !!! |%% 
%%%%%%%%%%%%%%%%%%%%%%%%%%%%%%%%%%%%%%%%%%%%%%%%%%%%%%%%%%%%%%%%%%%%%%%%%%%%%%%%%%%%%%%%%%%%%%%%%%%%%%%%%%%%%%%%%%%%%%%%%%%%%%%%%%%%%%%
%%| SET DOCUMENT PARAMTER HERE |%%

% !TEX root = Ausarbeitung.tex

\PassOptionsToPackage{noconfigs}{babel}
\documentclass[12pt]{article}

%%| INCLUDE PACKAGES HERE |%%
\usepackage[utf8]{inputenc}
\usepackage[paper=a4paper,left=25mm,right=25mm,top=25mm,bottom=25mm]{geometry}
\usepackage[english, ngerman]{babel}
\usepackage{graphicx}
\usepackage{amsmath}
\usepackage{amssymb}
\usepackage[autostyle=true,german=quotes]{csquotes}
\usepackage[backend=biber, style=ieee]{biblatex}
\usepackage{listings}
\usepackage{xcolor}

% 1. Load setspace BEFORE hyperref
\usepackage{setspace}

% 2. Load microtype (can be before or after setspace)
\usepackage{microtype}

% 3. Load hyperref LAST (standard rule for most LaTeX documents)
\usepackage{hyperref}

\usepackage{xurl} % Erlaubt Zeilenumbrüche an jeder Stelle in URLs

\usepackage{parskip} % Verhindert Einrücken, nutzt Abstände für Absätze


%%| SET PACKAGE PARAMETERS HERE |%%
\hypersetup{                                                                        
	colorlinks,
	citecolor   = black,
	filecolor   = black,
	linkcolor   = black,
	urlcolor    = black,
	pdftitle    = {Bericht},                                               % SET FILE PARAMETERS HERE (SHOWN IN PDF PROPERTIES)
	pdfsubject  = {Zwischenbericht Praxissemseter},                                       % ...
	pdfauthor   = {Clara Stoeckel},                                                 % ...
	pdfkeywords = {Zwischenbericht, Praxissemseter} ,                                    % ...
	pdfcreator  = {pdflatex},           
	pdfproducer = {LaTeX with hyperref}
}
\setstretch{1.25}                                                                   % SET GENERAL LINE SPACING HERE

%%| SET OTHER PARAMETERS HERE |%%
\addbibresource{ref.bib}


%%CodeExamples

% Farben definieren
\definecolor{codegreen}{rgb}{0,0.6,0}
\definecolor{codegray}{rgb}{0.5,0.5,0.5}
\definecolor{codepurple}{rgb}{0.58,0,0.82}
\definecolor{backcolour}{rgb}{0.95,0.95,0.92}

% Stil für den Code festlegen
\lstset{
	backgroundcolor=\color{backcolour},   % HIER: korrigiert von 'backroundcolor' zu 'backgroundcolor'
	commentstyle=\color{codegreen},
	keywordstyle=\color{magenta},
	numberstyle=\tiny\color{codegray},
	stringstyle=\color{codepurple},
	basicstyle=\ttfamily\footnotesize,
	breakatwhitespace=false,         
	breaklines=true,                 
	captionpos=b,                    
	keepspaces=true,                 
	numbers=left,                    
	numbersep=5pt,                  
	showspaces=false,                
	showstringspaces=false,
	showtabs=false,                  
	tabsize=2,
	frame=single
}

% QML als eigene Sprache definieren (da listings das nicht standardmäßig hat)
\lstdefinelanguage{QML}{
	keywords={import, property, function, id, target, anchors, Layout, qsTr, text, color, spacing, padding , title},
	sensitive=false,
	comment=[l]{//},
	morecomment=[s]{/*}{*/},
	morestring=[b]",
}


%%%%%%%%%%%%%%%%%%%%%%%%%%%%%%%%%%%%%%%%%%%%%%%%%%%%%%%%%%%%%%%%%%%%%%%%%%%%%%%%%%%%%%%%%%%%%%%%%%%%%%%%%%%%%%%%%%%%%%%%%%%%%%%%%%%%%%%
%%| TITLE PAGE |%%
\begin{document}
	\pagenumbering{Alph}
	\begin{titlepage}
		\begin{center}
			\includegraphics[width=0.55\textwidth]{HsKa_new.png}\\[14ex]
			\huge{\textbf{-Aktuelle Sicherheitsthemen-}\\Linux Kernel Exploit CVE-2024-1086}\\[9ex]
			%\large{\textbf{init}}\\[1ex]
			%\LARGE{\textbf{Telematic Devices (TDSW1)}}\\[8ex]
			%\normalsize{}
			%\large{\textbf{}}\\[1ex]
			%\LARGE{\textbf{}}\\[9ex]
			\normalsize{}
			\begin{tabular}{lll}
				Student:            & \quad Clara Stoeckel                                  & 93911 \\[2ex]     % NUMBER IS MAT.NR
				Universität:        & \quad Hochschule Karlsruhe – Technik und Wirtschaft   &       \\[2ex]
				Studiengang:        & \quad Informatik                &       \\[2ex]
				Semester:           & \quad Wintersemester 2025                             &       \\[2ex]
				Dozent:             & \quad Prof. Dr. Patrick Baier                       &       \\[2ex]
				Bearbeitet am:      & \quad 25. November 2025                                  &       \\[2ex]
			\end{tabular}
		\end{center}
	\end{titlepage}
	\newpage
%%%%%%%%%%%%%%%%%%%%%%%%%%%%%%%%%%%%%%%%%%%%%%%%%%%%%%%%%%%%%%%%%%%%%%%%%%%%%%%%%%%%%%%%%%%%%%%%%%%%%%%%%%%%%%%%%%%%%%%%%%%%%%%%%%%%%%%
%%| TABLE OF CONTENTS |%%
\pagenumbering{Roman}    % Intern eindeutige römische Zählung (I, II, ...)
\pagestyle{plain}        % Falls Seitenzahlen im Verzeichnis stehen sollen
% ODER:
% \pagenumbering{Roman} \pagestyle{empty} % Falls keine Zahlen im Verzeichnis stehen sollen

\tableofcontents
\newpage
	%%%%%%%%%%%%%%%%%%%%%%%%%%%%%%%%%%%%%%%%%%%%%%%%%%%%%%%%%%%%%%%%%%%%%%%%%%%%%%%%%%%%%%%%%%%%%%%%%%%%%%%%%%%%%%%%%%%%%%%%%%%%%%%%%%%%%%%
	%%| WRITE TEXT HERE |%%
	\pagenumbering{arabic}   % Startet die Zählung für den Hauptteil neu bei 1
	\setcounter{page}{1}     % Zur Sicherheit: Setzt den Zähler explizit auf 1
\section{Abstract}
Diese Ausarbeitung analysiert die Use-After-Free-Schwachstelle CVE-2024-1086 
im Netfilter-Subsystem \texttt{nf\_tables} des Linux-Kernels, welche alle 
Kernel-Versionen von 5.14 bis 6.6 betrifft und mit einem CVSS-Score von 7.8 
als hochkritisch eingestuft ist. Die Schwachstelle entsteht, weil die Funktion 
\texttt{nft\_verdict\_init()} positive Werte als Drop-Fehler im Hook-Urteil 
zulässt. Dies kann ein Double Free auslösen, wenn ein \texttt{NF\_DROP} mit 
einem \texttt{NF\_ACCEPT} ähnelnden Fehlercode ausgegeben wird~\cite{suse}. 
Ein lokaler Angreifer kann diesen Fehler durch gezielte Heap-Manipulation 
ausnutzen, um vollständige root-Rechte zu erlangen.

Ausgehend von den Grundlagen des Linux-Speichermanagements und der 
\texttt{nf\_tables}-Architektur wird der Fehler im Quellcode lokalisiert, der 
Exploit-Mechanismus nachvollzogen und der behebende Kernel-Patch analysiert. 
Abschließend werden bestehende Mitigationsmaßnahmen bewertet und ein Ausblick 
auf zukünftige Schutzmaßnahmen gegeben.
	\section{Grundlagen}
  \subsection{Der Linux Kernel \& Privilege Model}
  Der Linux Kernel bildet die Hauptkomponente des Linux Betriebssystem und dient als Schnittstelle zwischen Hardware und den Prozessen eines Computers.
  Hierbei hat er die folgenden vier grundlegende Funktionen: Speicherverwaltung, Prozessmanagement, Gerätetreiber und Systemaufrufe sowie Sicherheit. 
  Er agiert im Hintergrund und ist für den Nutzer unsichtbar, solange er korrekt implementiert ist. 
\cite{redhat}

  Eine sogenannte Userspace-Anwendung, also eine Anwendung die der Nutzer sieht und mit ihr agiert, interagiert mit dem Kernel über eine Systemaufrufschnittstelle (System Call Interface (SCI))
\cite{redhat}

Linux unterteilt die Berechtigungen des Root-Nutzers in seperate Privilegien (Capabilities). Privilegien können eine Teilmenge der Zugriffskontrollregeln des Unix-System umgehen.
Zum Beispiel die Berechtigungen eines Benutzers ändern.
Dabei verfügt jeder Linux-Prozess über drei Arten von Berechtigungen 

\begin{itemize}
	\item Effective (Effektiv): Berechtigungssatz der Entscheidungen über Zugriffskontrolle trifft
	\item Permitted (Zulässig): Liste von Privilegien, die der Prozess im aktiven Zustand verwenden darf
	\item Inheritable (Vererbbar): Berechtigungssatz, welche Privilegien ein Prozess erwerben kann
\end{itemize}
\cite{PrivAnalyzer}

Dies erklärt, wieso es fatal ist, wenn unberechtigte Personen Zugriff zu Rootrechten bekommen.
Sie können dann alle möglichen Änderungen auf Betriebssystemebene vornehmen, Daten lesen und schreiben, Nutzungsrechte verändern und mehr.

    % Ring-0 vs Ring-3, Kernel-/User-Space-Trennung
    % Capabilities (CAP_NET_ADMIN), User Namespaces
 \subsection{Speicherverwaltung im Kernel}

Der Kernel verwaltet seinen Speicher über die Funktionen \texttt{kmalloc()} 
und \texttt{kfree()}, sowie den SLUB-Allocator. Mittels \texttt{kmalloc()} 
kann Speicherplatz angefragt werden. Der SLUB-Allocator weist dann einen 
freien Slot in einem Speicher-Cache passender Größe zu. Objekte gleicher 
Größe werden dabei in gemeinsamen Pools, den sogenannten \textit{Slab-Caches}, 
verwaltet~\cite{love2010linux}. \texttt{kfree()} gibt den Speicherplatz 
wieder frei, sodass der SLUB-Allocator ihn erneut zuweisen kann. Der 
freigegebene Speicher wird dabei nicht sofort gelöscht oder gesperrt, 
sondern lediglich als verfügbar markiert~\cite{kernel_memory}.

Ein \textit{Use-After-Free} (UAF) entsteht, wenn Speicher freigegeben wurde, 
der Zeiger darauf jedoch weiterhin verwendet wird. Ein solcher Zeiger wird 
als \textit{dangling pointer} bezeichnet~\cite{love2010linux}. Da der 
SLUB-Allocator freigegebene Slots umgehend für neue Allokationen gleicher 
Größe wiederverwenden kann, besteht die Möglichkeit, dass ein anderes 
Kernel-Objekt denselben Speicherbereich belegt. Schreibt der dangling 
pointer anschließend in diesen Bereich, werden die Felder des neuen Objekts 
unkontrolliert überschrieben.

Ein \textit{Double Free} entsteht, wenn \texttt{kfree()} zweimal auf 
denselben Zeiger aufgerufen wird. Dabei wird die interne Verwaltungsstruktur 
des SLUB-Allocators korrumpiert, sodass derselbe Speicherbereich zwei 
verschiedenen Allokationen gleichzeitig zugewiesen werden kann~\cite{love2010linux}. 
Dies ist die direkte Ursache von CVE-2024-1086, wie in 
Abschnitt~\ref{sec:cve} beschrieben.

\textit{Heap Grooming} bezeichnet die gezielte Manipulation des Kernel-Heaps, 
um einen dangling pointer auf ein vom Angreifer kontrolliertes Objekt zu 
lenken~\cite{hawkes2008}. Das Vorgehen erfolgt in drei Schritten: Zunächst 
wird durch gezielte Allokationen der Heap in einen vorhersehbaren Zustand 
gebracht. Anschließend wird das Zielobjekt freigegeben. Zuletzt allokiert 
der Angreifer ein eigenes Objekt gleicher Größe, welches den freigewordenen 
Slot übernimmt. Der dangling pointer zeigt nun auf das kontrollierte Objekt.

Ein bevorzugtes Angriffsziel ist die \texttt{cred}-Struktur des Linux-Kernels. 
Sie speichert die sicherheitsrelevanten Attribute eines Prozesses, darunter 
die User-ID (\texttt{uid}) und Group-ID (\texttt{gid})~\cite{kernel_cred}. 
Gelingt es einem Angreifer, diese Felder auf den Wert \texttt{0} zu setzen, 
erhält der Prozess Root-Rechte. Genau dieses Prinzip liegt dem 
Exploit von CVE-2024-1086 zugrunde.
	\subsection{Netfilter \& nf\_tables}
    % Architektur: Hooks, Tables, Chains, Rules, Expressions
    % nft_verdict, nft_set – die relevanten Datenstrukturen

\section{CVE-2024-1086}
  \subsection{Verwundbare Komponente \& Ursache}
    % nft_verdict_init(): fehlende Validierung von
    % NF_DROP mit user-supplied verdict code
    % Zusammenspiel nft_chain_lookup_byid() → double free
  \subsection{Exploitation}
    \subsubsection{Heap Grooming \& SLUB-Manipulation}
      % Kontrolliertes Allokieren/Freigeben um
      % dangling pointer auf wiederverwendetes Objekt zu lenken
    \subsubsection{Privilege Escalation zur Root-Shell}
      % Überschreiben von cred-Strukturen
      % Notselwyns public PoC (write-up April 2024)
  \subsection{Patch \& Mitigationen}
    % Commit d0009ecc (fix in nft_verdict_init)
    % Gegenmaßnahmen: unprivileged user namespaces
    % deaktivieren, SMEP/SMAP, KASLR – Wirksamkeit

\section{Aktuelle Relevanz \& Einordnung}
  % CVSS 7.8, ausnutzbar seit Kernel 5.14 (2021!)
  % Verbreitung in Container-Umgebungen (Docker default config)
  % Zeitlinie: Discovery → PoC → Patches der Distributionen


	%%%%%%%%%%%%%%%%%%%%%%%%%%%%%%%%%%%%%%%%%%%%%%%%%%%%%%%%%%%%%%%%%%%%%%%%%%%%%%%%%%%%%%%%%%%%%%%%%%%%%%%%%%%%%%%%%%%%%%%%%%%%%%%%%%%%%%%
	%%| BIBLIOGRAPHY |%%                 
	\newpage                                          	
	
	% Print the Bibliography
	\printbibliography[heading=bibintoc, title={Literaturverzeichnis}]
	
	\newpage
	% List of Figures (Abbildungsverzeichnis)
	\addcontentsline{toc}{section}{Abbildungsverzeichnis}
	\listoffigures
	
	
	\end{document}